\documentclass[11pt]{article}
\usepackage[margin=1.2in]{geometry}
\usepackage{amsmath,amssymb,amsthm}
\usepackage{hyperref}
\hypersetup{colorlinks=true, linkcolor=blue, urlcolor=blue}
\usepackage{parskip}
\emergencystretch=3em

\newcommand{\R}{\mathbb{R}}
\newcommand{\code}[1]{\texttt{#1}}

\title{Tide retrospective: singular composition, part 1\\
\large the critique's prep lemmas}
\author{Claude (Anthropic), with GPT-5.6 Sol consults}
\date{2026-08-10}

\begin{document}
\maketitle

\begin{abstract}
An adversarial GPT-5.6 Sol audit of the note's unnormalized
degenerate theorem (\code{thm:singular}) --- commissioned to find
holes in the informal and formal proofs, and archived verbatim in the
tide log --- found no mathematical flaw in either, but two moderate
coverage findings: the Lean corpus verifies the quantitative core
(pencil, sector, lower bound, contradiction) rather than the
note-level statement, with twelve routine reduction steps left
semantic; and the note's footnote overclaimed the machine-checked
fragment. On user approval, the footnote and the two integer-exponent
phrasings were fixed in the TeX, and this tide begins closing the
composition gap with the three hypothesis-manufacturing lemmas: the
quadratic upper bound $K \le C_0\|w\|^2$ at a nonnegative $C^2$ zero
(gradient by the local-minimum rule, second order by the Lipschitz
derivative and the mean value inequality); the least nonzero diagonal
of a nonvanishing analytic function, with a witness rescaled to the
endpoint's $\|x_0\| = 3/2$ normalization --- and here the critique's
own plan simplified, since the power-series sum evaluates only
diagonals, so no symmetric-multilinear polarization is needed; and a
hand-rolled continuous bump equal to one on a prescribed ball.
\end{abstract}

\section{Setting}

The critique's section B4 enumerates precisely what stands between
the corpus and Theorem \code{thm:singular}: point selection,
translation, power-series extraction, least degree and direction,
the quadratic bound, the bump, the local-to-global reduction, the
expansion-family bridge, the integral rewriting, and the neighborhood
assembly. This tide delivers the middle block (steps 4--7); the
composed point theorem and the locus assembly follow as parts 2--3.

\section{The thing formalised}

\begin{sloppypar}
\code{quadratic\_upper\_bound\_of\_nonneg}:
\code{IsLocalMin.fderiv\_eq\_zero} kills the gradient;
\code{ContDiffAt.fderiv\_right} plus
\code{exists\_lipschitzOnWith} make the derivative $M$-Lipschitz
near the zero; the \code{hasFDerivWithin} mean-value bound on the
$\|w\|$-ball then gives $K(w) \le M\|w\|^2$ --- the catalogued
unification-bomb warning steered the proof to the
explicit-derivative form from the start.
\code{exists\_least\_nonzero\_diagonal}: the constant diagonal is
killed first (\code{hasSum\_single} at the center plus
\code{map\_coord\_zero}), so \code{Nat.find} lands at degree at
least one without any dependent-index transport; homogeneity
(\code{map\_smul\_univ}) rescales the witness.
\code{exists\_bump\_one\_on\_ball}: $\min(1, \max(0, R+1-\|w\|))$ ---
continuous, supported in the $(R{+}1)$-ball, one on the $R$-ball;
the endpoint requires only continuity, so smooth-bump instances
(which would drag in norm-smoothness questions on the sup-normed
pi space) are avoided entirely.
\end{sloppypar}

\section{Where progress stalled}

\begin{sloppypar}
Three rounds. \code{ContDiffAt.differentiableAt} now takes
$n \ne 0$; an \code{interval\_cases} on a \code{set}-bound
\code{Nat.find} index tangles with the dependent \code{Fin} type in
the multilinear argument (the restructure --- prove the degenerate
case's fact generally first, then rewrite inside the \code{find}
spec --- is the reusable move); and \code{EMetric.ball} is deprecated
for \code{Metric.eball} on this pin.
\end{sloppypar}

\section{Mathlib lookup versus new mathematics}

\begin{sloppypar}
Lookup: the local-minimum derivative rule
(\code{IsLocalMin.fderiv\_eq\_zero}), the local Lipschitz bound
(\code{exists\_lipschitzOnWith}), the convex mean-value inequality
(\code{norm\_image\_sub\_le\_of\_norm\_hasFDerivWithin\_le}), the
power-series sum (\code{HasFPowerSeriesOnBall.hasSum}),
\code{hasSum\_single}, \code{map\_coord\_zero},
\code{map\_smul\_univ}, and \code{Nat.find\_min}. New: only the
observation that diagonal-only summation makes polarization
unnecessary --- a simplification over both the note's argument and
the critique's own reconstruction.
\end{sloppypar}

\section{What the next tide inherits}

Part 2: the composed point theorem. Translation of the pencil
integrals to an arbitrary center (Lebesgue substitution), the
expansion-family bridge (equality of families at the observable
$g\psi$ yields the $o(t^{-N})$-for-every-$N$ premise) and the
integral rewriting (critique steps 10--11), then
\code{pencil\_families\_force\_germ\_eq\_at} by composing this
tide's three lemmas with
\code{analytic\_pencil\_difference\_not\_superpolynomial}. Part 3
assembles the locus statement by compactness. The TeX footnote will
be updated once more when the composition lands, to say what is then
true: the theorem is machine-checked end to end.
\end{document}
